Análisis
La evolución mensual de los saldos bancarios durante 2025 evidencia un comportamiento caracterizado por fluctuaciones a lo largo del período analizado. El saldo total administrado por el banco oscila entre aproximadamente
Q1.79 millones
y
Q2.55 millones
, registrándose el valor más alto en el mes de diciembre, lo que representa una variación cercana a
Q800 mil
entre el valor mínimo y el máximo observado.
No se identifica una tendencia sostenida de crecimiento o disminución durante el año. En su lugar, los saldos presentan variaciones mensuales que podrían estar asociadas a factores estacionales o a cambios en el comportamiento financiero de los clientes.
El incremento registrado en diciembre podría estar relacionado con eventos propios de la época, como el pago de aguinaldos, bonos, remesas o un mayor dinamismo de la actividad económica. No obstante, confirmar estas hipótesis requeriría incorporar información adicional, como indicadores macroeconómicos o campañas comerciales desarrolladas por la institución durante el mismo período.